% Generated draft: requires mathematical and source review.
\documentclass[11pt]{article}
\usepackage[margin=25mm]{geometry}
\usepackage{fontspec}
\setmainfont{Latin Modern Roman}
\usepackage{amsmath,amssymb,mathtools,unicode-math}
\setmathfont{Latin Modern Math}
\usepackage{enumitem,needspace,xurl,hyperref}
\setlength{\parindent}{0pt}
\setlength{\parskip}{5pt}
\setlength{\emergencystretch}{3em}
\newcommand{\N}{\mathbb{N}}
\newcommand{\Z}{\mathbb{Z}}
\newcommand{\Q}{\mathbb{Q}}
\newcommand{\R}{\mathbb{R}}
\newcommand{\C}{\mathbb{C}}
\newcommand{\F}{\mathbb{F}}
\newcommand{\A}{\mathbb{A}}
\newcommand{\PP}{\mathbb{P}}
\newcommand{\E}{\mathbb{E}}
\newcommand{\eps}{\varepsilon}
\newcommand{\dd}{\,\mathrm{d}}
\begin{document}
\section*{UnsolvedMath ID: 701388}
\textbf{Source problem code:} catalogue entry
\subsection{Definitions and mathematical statement}
Let \(F=F(x_1,x_2,x_3)\) be the free group on three generators. For \(r_1,r_2,r_3\in F\), assume that
\[
F/\langle\!\langle r_1,r_2,r_3\rangle\!\rangle_F=\{1\},
\]
where \(\langle\!\langle S\rangle\!\rangle_F\) denotes the smallest normal subgroup of \(F\) containing \(S\). Thus \(\langle x_1,x_2,x_3\mid r_1,r_2,r_3\rangle\) presents the trivial group. The source describes the following as simple questions related to the Whitehead asphericity conjecture.

\begin{enumerate}
\item[(a)] For every triple \((r_1,r_2,r_3)\) satisfying the assumption above, prove that
\[
G=F/\langle\!\langle r_1,r_2\rangle\!\rangle_F
\]
is torsion-free. Explicitly, the assertion is
\[
\forall g\in G\ \forall n\in\Z_{\geq 1},\qquad g^n=1\Longrightarrow g=1.
\]
\item[(b)] Put \(R_i=\langle\!\langle r_i\rangle\!\rangle_F\) for \(i=1,2,3\). For subgroups \(A,B\leq F\), write
\[
[A,B]=\langle a^{-1}b^{-1}ab\mid a\in A,\ b\in B\rangle.
\]
Since \(R_1\) and \(R_2\) are normal in \(F\), their commutator subgroup \([R_1,R_2]\) is normal in \(F\). The archived question asks whether, for every triple satisfying the same assumption,
\[
H=F/[R_1,R_2]
\]
is residually soluble. Here a group \(S\) is soluble if its derived series, defined by \(S^{(0)}=S\) and \(S^{(k+1)}=[S^{(k)},S^{(k)}]\), reaches \(\{1\}\) after finitely many steps. A group \(H\) is residually soluble if, for every \(h\in H\setminus\{1\}\), there exist a soluble group \(S\) and a homomorphism \(\varphi:H\to S\) with \(\varphi(h)\ne 1\). No finiteness condition on \(S\) is imposed.
\end{enumerate}

The supplied record places part (a) in the remaining main-section material and part (b) in the Archive of Solved Problems. It reports a negative answer to part (b), attributed to V. A. Roman'kov. These are supplied source and status qualifications, not independently verified conclusions of this draft; no counterexample or proof is supplied here.
\subsection{Short English statement}
Starting with any three-generator, three-relator presentation of the trivial group, part (a) asks whether deleting the third relator always leaves a torsion-free group. Part (b) asks whether quotienting the free group by the commutator of the normal closures of the first two relators always gives a residually soluble group. The supplied record reports part (b) as negatively resolved and part (a) as remaining.
\subsection{Sources}
\begin{itemize}
\item Supplied local UnsolvedMath registry record: permanent ID \(701388\), title \emph{Kourovka Notebook Problem 17.86}, exact source problem code \texttt{catalogue entry}; upstream revision \texttt{d371ee917ac9647fbb06cd6fc529deec6f3dda82}.
\item \emph{The Kourovka Notebook (Unsolved Problems in Group Theory)}, 17th issue (2010), Problem 17.86, proposed by R. Mikhailov. The supplied background identifies an update in the 21st issue (2026), \texttt{arXiv:1401.0300v48}. The supplied collection description names E. I. Khukhro and V. D. Mazurov as editors.
\item Supplied Archive of Solved Problems entry for part (b): negative answer attributed to V. A. Roman'kov, \emph{Siberian Math. J.}, volume 52, number 2 (2011), pages 348--351. The supplied record gives no article title or direct article link.
\item The supplied status review, dated exactly \texttt{2026-10-16}, labels the problem \texttt{partially\_solved}, records its evidence as published, and states that part (b) is archived as solved while the remaining part is in the main section. This assessment is supplied dated curation, not independently verified literature or source verification.
\end{itemize}
\end{document}
